\documentclass[10pt,a4paper]{article}

\usepackage[margin=1.45cm,top=1.25cm,bottom=1.3cm]{geometry}
\usepackage[T1]{fontenc}
\usepackage[utf8]{inputenc}
\usepackage{lmodern}
\usepackage{microtype}
\usepackage{xcolor}
\usepackage{enumitem}
\usepackage{titlesec}
\usepackage[hidelinks]{hyperref}
\usepackage{tabularx}
\usepackage{array}
\usepackage{lastpage}
\usepackage{fancyhdr}
\input{glyphtounicode}
\pdfgentounicode=1
\hypersetup{
  pdftitle={Divyesh A. Mistry - Computational Materials Scientist},
  pdfauthor={Dr. Divyesh A. Mistry},
  pdfsubject={Curriculum Vitae},
  pdfkeywords={computational materials scientist, molecular dynamics, LAMMPS, discrete dislocation dynamics, multiscale modeling, Python, high-performance computing, fracture mechanics, plasticity, scientific software}
}

\definecolor{navy}{HTML}{17365D}
\definecolor{rulegray}{HTML}{A7B0BA}
\definecolor{textgray}{HTML}{333333}
\color{textgray}

\setlength{\parindent}{0pt}
\setlength{\parskip}{0pt}
\setlength{\emergencystretch}{2em}
\setlist[itemize]{leftmargin=1.15em,label=--,itemsep=1.2pt,topsep=2pt,parsep=0pt}
\titleformat{\section}{\large\bfseries\color{navy}}{}{0pt}{}[\vspace{-2pt}\color{rulegray}\titlerule]
\titlespacing*{\section}{0pt}{7pt}{4pt}
\pagestyle{fancy}
\fancyhf{}
\renewcommand{\headrulewidth}{0pt}
\fancyfoot[C]{\footnotesize Divyesh Mistry | Curriculum Vitae | Page \thepage\ of \pageref{LastPage}}

\newcommand{\role}[4]{%
  \textbf{#1} \hfill \textbf{#2}\\[-1pt]
  \textit{#3} \hfill \textit{#4}\vspace{1pt}
}

\begin{document}

{\LARGE\bfseries\color{navy} Divyesh A. Mistry, PhD}\hfill
{\large\bfseries Computational Materials Scientist}\par
\vspace{3pt}
Seville, Spain \;|\; +34 694 290 182 \;|\;
\href{mailto:d.mistryg@gmail.com}{d.mistryg@gmail.com} \;|\;
\href{https://www.linkedin.com/in/dr-divyesh-mistry-7b40a95b/}{linkedin.com/in/dr-divyesh-mistry} \;|\;
\href{https://github.com/divyesh-mistry}{github.com/divyesh-mistry} \;|\;
\href{https://scholar.google.com/citations?user=feQuDO8AAAAJ\&hl=en}{Google Scholar}

\section{Professional Summary}
Computational materials scientist specializing in molecular dynamics, multiscale mechanics, and high-performance scientific computing. Develops LAMMPS, discrete dislocation dynamics, and Python workflows that connect atomistic mechanisms to measurable deformation, strengthening, adhesion, and fracture response. Experience spans Ni-based superalloys, tungsten and W-Re alloys, polymer adhesives, constitutive modeling, defect characterization, and reproducible analysis on HPC clusters. Published first-author research and open simulation software; available for computational materials, scientific software, and mechanics roles.

\section{Technical Skills}
\textbf{Simulation \& Modeling:} Molecular dynamics (MD), discrete dislocation dynamics (DDD), finite element analysis (FEA/FEM), multiscale modeling, constitutive modeling, fracture mechanics, plasticity, microstructure evolution.\\[2pt]
\textbf{Simulation Software:} LAMMPS, OVITO, Atomsk, ABAQUS, ANSYS, COMSOL, NEPER, ParaView, MATLAB.\\[2pt]
\textbf{Programming \& Scientific Python:} Python (NumPy, Pandas, SciPy, Matplotlib), C++, Fortran, Bash, workflow automation, large-scale data analysis, Git/GitHub, \LaTeX.\\[2pt]
\textbf{High-Performance Computing (HPC):} Linux/Unix, MPI/OpenMP parallel computing, HPC cluster management, batch scheduling, multi-million-atom simulations.

\section{Professional Experience}
\role{Postdoctoral Researcher}{Aug 2026--Present}{University of Seville}{Seville, Spain}
\begin{itemize}
  \item Develop molecular-dynamics and multiscale workflows for cross-linked epoxy and polymer adhesive systems, linking network chemistry and crosslink density to thermomechanical response, glass-transition temperature, shear strength, and failure.
  \item Design reproducible LAMMPS protocols for curing, equilibration, cooling, and deformation; evaluate thermostat, strain-rate, and finite-size effects to improve verification and physical interpretation.
  \item Build Python pipelines for replica-based statistical analysis, stress--strain processing, property extraction, uncertainty assessment, and automated visualization of large simulation campaigns on HPC systems.
  \item Connect atomistic shear response with cohesive-zone and continuum descriptions of adhesive joints to support strength and life-prediction models across length scales.
  \item \textbf{Project:} \textit{Multiscale modeling of interfacial delamination of steel/adhesive interfaces by aqueous corrosion}, funded by the U.S. Office of Naval Research.
  \item \textbf{Advisors:} Prof. Mar\'{i}a del Pilar Ariza Moreno and Prof. Michael Ortiz.
\end{itemize}

\role{Visiting Researcher}{Apr 2026--Jul 2026}{Merrimack College, Department of Mechanical and Electrical Engineering}{North Andover, MA, USA}
\begin{itemize}
  \item Performed large-scale MD simulations of temperature-dependent crack-tip deformation and the ductile-to-brittle transition in tungsten and W--Re alloys.
  \item Developed automated LAMMPS and Python workflows for stress--strain analysis, fracture-energy calculations, defect characterization, dislocation extraction, and trajectory processing on HPC clusters.
  \item Integrated atomistic fracture mechanisms with a 3D polycrystalline DDD framework to study Frank--Read source activation, dislocation multiplication, grain-boundary interactions, and plastic shielding.
  \item Mentored student research and contributed to manuscript development and multiscale simulation planning.
\end{itemize}

\role{Junior Research Fellow, DRDO-Sponsored Project}{Jan 2025--Dec 2025}{Indian Institute of Technology Bombay}{Mumbai, India \normalfont\textit{(concurrent with final thesis defense)}}
\begin{itemize}
  \item Investigated constitutive behavior and high-strain-rate response of shape-memory-alloy composites for aerospace and defense applications.
  \item Conducted Split Hopkinson Pressure Bar experiments and used experimental data to calibrate and validate numerical constitutive models.
  \item Supported project execution, data analysis, technical reporting, and mentoring of junior researchers.
\end{itemize}

\role{PhD Researcher, Aerospace Engineering}{Jan 2019--Jun 2025}{Indian Institute of Technology Bombay}{Mumbai, India}
\begin{itemize}
  \item Developed an atomistically informed multiscale framework for dislocation--precipitate--interface interactions in powder-metallurgy Ni-based superalloys.
  \item Executed multi-million-atom LAMMPS simulations to quantify generalized stacking-fault energies, dislocation mobility, precipitate shearing, critical resolved shear stress, and temperature-dependent mechanisms.
  \item Established size-dependent power laws for paired edge dislocations interacting with $\gamma'$ precipitates and transferred atomistic drag and obstacle-strength parameters into DDD models.
  \item Developed polycrystalline DDD/finite-element models of prior particle boundaries, quantifying hardening, pinned dislocation density, geometrically necessary dislocations, and residual-stress retention.
  \item Created parallel and automated Python/C++/Fortran workflows for simulation setup, defect tracking, post-processing, verification, parameter studies, and publication-quality visualization.
  \item Published first-author research, presented at international conferences, and collaborated with computational and experimental researchers.
\end{itemize}

\role{Graduate Teaching Assistant}{Jan 2019--Dec 2023}{Indian Institute of Technology Bombay}{Mumbai, India}
\begin{itemize}
  \item Supported Continuum Mechanics, Data Analysis, and Multiscale Modeling courses; developed computational exercises and guided undergraduate and postgraduate students.
\end{itemize}

\role{Assistant Professor, Mechanical Engineering}{Jun 2016--Dec 2018}{CMR Institute of Technology}{Bengaluru, India}
\begin{itemize}
  \item Taught Finite Element Methods, Experimental Stress Analysis, and Solid Mechanics; developed laboratory modules using ANSYS, MATLAB, and Python.
\end{itemize}

\role{Graduate Teaching Assistant and MTech Researcher}{Jun 2014--Jul 2016}{Indian Institute of Technology Kharagpur}{Kharagpur, India}
\begin{itemize}
  \item Developed and validated ABAQUS and ANSYS models for structural mechanics, smart composites, and vibration; supported laboratories in FEM, vibration analysis, and tribology.
\end{itemize}

\section{Education}
\textbf{PhD, Aerospace Engineering}, Indian Institute of Technology Bombay \hfill 2019--2025\\
Thesis: \textit{Multiscale Modeling of Prior Particle Boundaries in Nickel-Based Superalloys}; CGPA: 9.36/10\\
Advisors: Prof. P. J. Guruprasad and Prof. Amuthan A. Ramabathiran

\vspace{3pt}
\textbf{MTech, Mechanical Engineering (Mechanical System Design)}, IIT Kharagpur \hfill 2014--2016\\
Thesis: \textit{A Mesh-Free Model for Static Analysis of Smart Composite Beams}; CGPA: 8.85/10

\vspace{3pt}
\textbf{BE, Aeronautical Engineering (Aero Mechanical)}, Aeronautical Society of India \hfill 2013

\section{Publications}
\textbf{\textit{Peer-Reviewed Journal Articles}}
\begin{enumerate}[leftmargin=1.35em,itemsep=2pt,topsep=1pt]
  \item \textbf{Mistry, D. A.}, Tak, T. N., Sankarasubramanian, R., and Guruprasad, P. J. ``A Discrete Dislocation Dynamics Study of Prior Particle Boundary Effects in Ni-Based Superalloys.'' \textit{Journal of Engineering Materials and Technology} (2026). \href{https://doi.org/10.1115/1.4072448}{DOI}.
  \item \textbf{Mistry, D. A.} and Ramabathiran, A. A. ``Size-Dependent Power Laws for Edge Dislocations in Nickel Superalloys: A Molecular Dynamics Study.'' \textit{Computational Materials Science} 259, 114122 (2025). \href{https://doi.org/10.1016/j.commatsci.2025.114122}{DOI}.
\end{enumerate}

\textbf{\textit{Manuscripts in Preparation}}
\begin{enumerate}[leftmargin=1.35em,itemsep=2pt,topsep=1pt]
  \item \textbf{Mistry, D. A.} and Mahata, A. ``Atomistic Indicators of the Ductile-to-Brittle Transition in Polycrystalline Tungsten: Temperature and Rhenium Effects on Crack-Tip Plasticity.'' arXiv preprint arXiv:2608.27589 (2026). \href{https://arxiv.org/abs/2608.27589}{Preprint}.
  \item \textbf{Mistry, D. A.}, Tak, T. N., and Guruprasad, P. J. ``Polycrystalline Discrete Dislocation Dynamics of Ni-Based Superalloys: Interactions with Prior Particle Boundaries and Second Phases Using Atomistically Informed Inputs.'' Manuscript in preparation.
\end{enumerate}

\textbf{\textit{Conference Contributions}}
\begin{enumerate}[leftmargin=1.35em,itemsep=2pt,topsep=1pt]
  \item EMI 2026, Boulder, USA: \textit{Multiscale Modeling of Dislocation--Precipitate--PPB Interactions in Powder-Metallurgy Nickel Superalloys}.
  \item IMPLAST 2025, IIT Roorkee, India: \textit{Hierarchical Multiscale Modeling of Plasticity in Ni-Based Superalloys}.
  \item COMPLAS 2023, Barcelona, Spain: \textit{An Atomistically Informed DDD Study of Prior Particle Boundaries in Ni-Based Superalloys}.
\end{enumerate}


\section{Awards, Grants \& Professional Memberships}
\begin{itemize}
  \item \textbf{Institution of Eminence Travel Grant}, IIT Bombay (2025) --- funded attendance at TMS 2025, USA.
  \item \textbf{IIT Bombay International Travel Grant} (2023) --- funded attendance at COMPLAS 2023, Spain.
  \item \textbf{GATE All India Rank: 69} (2014) --- Graduate Aptitude Test in Engineering; MHRD fellowships for doctoral and master's study.
  \item \textbf{Member}, The Minerals, Metals \& Materials Society (TMS); \textbf{Associate Member}, Aeronautical Society of India (AeSI).
  \item Certifications: Using Python for Research (HarvardX); The Unix Workbench (Johns Hopkins); Material Behavior (Georgia Tech); Engineering Simulations (CornellX).
\end{itemize}

\section{Additional Information}
Research interests: computational materials science, multiscale mechanics, atomistic simulation, deformation and fracture, additive-manufacturing materials, polymer adhesives, and data-driven materials modeling. Academic references available upon request.

\end{document}
